\textbf{Phần 1. Trắc nghiệm lựa chọn}
\setcounterpageref{ex}{0}

\begin{ex} %[4H6V29]
Hai đường thẳng song song là hai đường thẳng
\choice
{Cắt nhau tại 1 điểm}
{\True Không bao giờ cắt nhau (dù kéo dài mãi)}
{Vuông góc với nhau}
{Trùng nhau}
\loigiai{
Hai đường thẳng song song không có điểm chung nào.\\
\textbf{Đáp án của câu hỏi là B}
}
\end{ex}

\begin{ex} %[4H6V29]
Hình nào sau đây thể hiện hình ảnh của hai đường thẳng song song trong thực tế?
\choice
{Hai kim đồng hồ lúc 3 giờ}
{Hai lưỡi kéo đang mở}
{Hai cạnh của góc vuông ê-ke}
{\True Hai thanh ray của đường tàu hỏa}
\loigiai{
Hai thanh ray đường tàu hỏa song song với nhau.\\
\textbf{Đáp án của câu hỏi là D}
}
\end{ex}

\begin{ex} %[4H6V29]
Trong hình chữ nhật $ABCD$, cặp cạnh nào sau đây song song với nhau?
\choice
{\True $AB$ và $CD$}
{$AB$ và $BC$}
{$BC$ và $CD$}
{$AB$ và $AD$}
\loigiai{
Hai cạnh đối diện $AB$ và $CD$ song song với nhau.\\
\textbf{Đáp án của câu hỏi là A}
}
\end{ex}

\begin{ex} %[4H6V29]
Hình bình hành có mấy cặp cạnh đối diện song song và bằng nhau?
\choice
{$1$ cặp}
{$3$ cặp}
{\True $2$ cặp}
{$4$ cặp}
\loigiai{
Hình bình hành có 2 cặp cạnh đối diện song song và bằng nhau.\\
\textbf{Đáp án của câu hỏi là C}
}
\end{ex}

\begin{ex} %[4H6V29]
Hình thang là hình tứ giác có
\choice
{2 cặp cạnh đối diện song song}
{\True 1 cặp cạnh đối diện song song (hai đáy)}
{4 cạnh bằng nhau}
{4 góc vuông}
\loigiai{
Hình thang có 1 cặp cạnh đối diện song song.\\
\textbf{Đáp án của câu hỏi là B}
}
\end{ex}

\begin{ex} %[4H6V29]
Trong hình vuông $ABCD$, cạnh $AD$ song song với cạnh nào?
\choice
{$AB$}
{$CD$}
{$AC$}
{\True $BC$}
\loigiai{
Cạnh $AD$ song song với cạnh đối diện $BC$.\\
\textbf{Đáp án của câu hỏi là D}
}
\end{ex}

\begin{ex} %[4H6V29]
Hai đường thẳng cùng vuông góc với một đường thẳng thứ ba thì chúng
\choice
{\True Song song với nhau}
{Cắt nhau}
{Vuông góc với nhau}
{Trùng nhau}
\loigiai{
Nếu $a \perp c$ và $b \perp c$ thì $a // b$.\\
\textbf{Đáp án của câu hỏi là A}
}
\end{ex}

\begin{ex} %[4H6V29]
Qua một điểm $M$ nằm ngoài đường thẳng $a$, có thể vẽ được bao nhiêu đường thẳng song song với $a$?
\choice
{$0$}
{$2$}
{\True $1$ đường thẳng duy nhất}
{Vô số}
\loigiai{
Qua 1 điểm ngoài đường thẳng chỉ vẽ được 1 đường thẳng song song duy nhất.\\
\textbf{Đáp án của câu hỏi là C}
}
\end{ex}

\begin{ex} %[4H6V29]
Khoảng cách giữa hai đường thẳng song song luôn
\choice
{Thay đổi}
{\True Không đổi (bằng nhau tại mọi điểm)}
{Bằng 0}
{Tăng dần}
\loigiai{
Khoảng cách giữa hai đường thẳng song song không đổi.\\
\textbf{Đáp án của câu hỏi là B}
}
\end{ex}

\begin{ex} %[4H6V29]
Trong các hình sau, hình nào KHÔNG có cặp cạnh nào song song?
\choice
{Hình chữ nhật}
{Hình bình hành}
{Hình thoi}
{\True Hình tam giác}
\loigiai{
Hình tam giác có 3 cạnh cắt nhau từng đôi một, không có cạnh song song.\\
\textbf{Đáp án của câu hỏi là D}
}
\end{ex}

\begin{ex} %[4H6V29]
Hai dòng kẻ ngang trên trang vở ô li là hình ảnh của
\choice
{\True Hai đường thẳng song song}
{Hai đường thẳng vuông góc}
{Hai đường thẳng cắt nhau}
{Hai tia chung gốc}
\loigiai{
Các dòng kẻ ngang song song với nhau.\\
\textbf{Đáp án của câu hỏi là A}
}
\end{ex}

\begin{ex} %[4H6V29]
Cho hình thoi $ABCD$. Cặp cạnh nào sau đây song song với nhau?
\choice
{$AB$ và $BC$}
{$AC$ và $BD$}
{\True $AB$ và $CD$ (cũng như $AD$ và $BC$)}
{$AD$ và $DC$}
\loigiai{
Hình thoi có các cặp cạnh đối diện song song: $AB // CD$ và $AD // BC$.\\
\textbf{Đáp án của câu hỏi là C}
}
\end{ex}

\textbf{Phần 2. Trắc nghiệm đúng sai}
\setcounterpageref{ex}{0}

\begin{ex} %[4H6V29]
Xét tính đúng sai về hai đường thẳng song song:
\choiceTF
{\True Hai đường thẳng song song không bao giờ cắt nhau.}
{\True Trong hình chữ nhật có 2 cặp cạnh đối diện song song.}
{Hai đường thẳng không cắt nhau thì luôn vuông góc.}
{\True Hai mép đối diện của cánh cửa lớp học là hình ảnh của hai đoạn thẳng song song.}
\loigiai{
\begin{itemize}
    \item Song song không cắt nhau: Đúng.
    \item Hình chữ nhật có 2 cặp cạnh song song: Đúng.
    \item Không cắt nhau là song song (không phải vuông góc). Sai.
    \item Hai mép cánh cửa song song: Đúng.
\end{itemize}
\begin{itemize}
    \item a) \textbf{Đúng}.
    \item b) \textbf{Đúng}.
    \item c) \textbf{Sai}.
    \item d) \textbf{Đúng}.
\end{itemize}
}
\end{ex}

\begin{ex} %[4H6V29]
Cho hình bình hành $ABCD$. Xét tính đúng sai:
\choiceTF
{\True Cạnh $AB$ song song với cạnh $CD$.}
{\True Cạnh $AD$ song song với cạnh $BC$.}
{Cạnh $AB$ song song với cạnh $BC$.}
{\True Độ dài cạnh $AB$ bằng độ dài cạnh $CD$.}
\loigiai{
\begin{itemize}
    \item $AB // CD$: Đúng.
    \item $AD // BC$: Đúng.
    \item $AB$ cắt $BC$ tại $B$ (không song song). Sai.
    \item $AB = CD$: Đúng.
\end{itemize}
\begin{itemize}
    \item a) \textbf{Đúng}.
    \item b) \textbf{Đúng}.
    \item c) \textbf{Sai}.
    \item d) \textbf{Đúng}.
\end{itemize}
}
\end{ex}

\textbf{Phần 3. Trả lời ngắn}
\setcounterpageref{ex}{0}

\begin{ex} %[4H6V29]
Hình chữ nhật có bao nhiêu cặp cạnh đối diện song song?
\shortans{2}
\loigiai{
\textbf{Đáp số:} $2$\\
2 cặp cạnh song song.
}
\end{ex}

\begin{ex} %[4H6V29]
Hình vuông có bao nhiêu cặp cạnh đối diện song song?
\shortans{2}
\loigiai{
\textbf{Đáp số:} $2$\\
2 cặp cạnh song song.
}
\end{ex}

\begin{ex} %[4H6V29]
Hình thang thông thường có bao nhiêu cặp cạnh đối diện song song?
\shortans{1}
\loigiai{
\textbf{Đáp số:} $1$\\
1 cặp cạnh song song (hai đáy).
}
\end{ex}

\begin{ex} %[4H6V29]
Qua một điểm $M$ ngoài đường thẳng $d$, vẽ được bao nhiêu đường thẳng song song với $d$?
\shortans{1}
\loigiai{
\textbf{Đáp số:} $1$\\
1 đường thẳng duy nhất.
}
\end{ex}

\begin{ex} %[4H6V29]
Hình bình hành có bao nhiêu cặp cạnh đối diện song song?
\shortans{2}
\loigiai{
\textbf{Đáp số:} $2$\\
2 cặp.
}
\end{ex}

\begin{ex} %[4H6V29]
Một hình lục giác đều có bao nhiêu cặp cạnh đối diện song song với nhau?
\shortans{3}
\loigiai{
\textbf{Đáp số:} $3$\\
3 cặp cạnh đối diện song song.
}
\end{ex}
